%% ============================================================================================
%% Port numbering specification for local clouds and systems of local clouds.
%% Normative section of the Core GSoSDD; referenced by the Support Systems GSoSDD.
%% ============================================================================================
\section{Port Numbering Specification}
\label{sec:ports}

Arrowhead systems discover each other through the \sys{ServiceRegistry}, so in principle any
port works. In practice fixed, well-known default ports matter: the \sys{ServiceRegistry} port
must be known for bootstrapping, test and reference deployments put many systems on one host,
firewall rules and container mappings are written against them, and several local clouds are
often run side by side. Up to 5.2.2 the defaults were assigned per repository without a common
plan, which led to the same default port (8465) being used by
\sys{SimpleStoreServiceOrchestration}, \sys{TranslationManager} and \sys{Onboarding}. This
section specifies how ports are allocated within a local cloud and across a system of local
clouds. The key words \emph{shall}, \emph{should} and \emph{may} are used as in RFC 2119.

\subsection{Principles}

\begin{description}
  \item[P1 -- One canonical port per system type.] Every Arrowhead framework system type (Core
        and Support) \emph{shall} have exactly one canonical default port, unique across the
        whole framework allocation (Table~\ref{tab:portalloc}). Any combination of framework
        systems can thereby be co-located on one host without reconfiguration.
  \item[P2 -- One port per system instance.] A system \emph{shall} serve all its HTTP(S)
        services on its canonical port; the same port is used in secure and insecure mode.
        Additional listeners (health or sidecar endpoints) \emph{shall} use the health range.
  \item[P3 -- Stability.] Allocated ports are never changed without a Roadmap WG decision. Ports
        of retired systems stay \emph{reserved} for at least one major version and are not
        re-assigned. The Arrowhead 4.6 allocations are kept wherever they are unique.
  \item[P4 -- Discovery over configuration.] Only the \sys{ServiceRegistry} port is needed for
        bootstrapping. Consumers \emph{shall not} hard-code any other port; the registered
        \prop{accessPort} of a service interface is authoritative.
  \item[P5 -- Reachability of registered ports.] The port a system registers \emph{shall} be
        the port its consumers use. If address or port translation is applied (NAT, container
        port publishing), the registered value \emph{shall} be valid for the consumers of the
        local cloud, or translation \emph{shall} follow the cloud offset rule of
        Section~\ref{sec:ports:sos}.
  \item[P6 -- Ephemeral-range avoidance.] Allocated ports \emph{shall} lie below 32768, the
        start of the Linux ephemeral port range, so that outgoing connections never occupy them.
\end{description}

\subsection{Port Block of a Local Cloud}
\label{sec:ports:lc}

A local cloud uses the 1000-port \emph{block} 8000--8999 (cloud index $k=0$), divided into
functional ranges (Table~\ref{tab:portranges}).

\begin{table}[h!]
\centering\footnotesize
\begin{tabularx}{\textwidth}{|P{2.2cm}|P{5.0cm}|Y|}\hline
\rowcolor{gray!33}Range & Use & Allocation authority \\\hline
8000--8099 & Tunnel control port pool (\sys{GateTunnel} session control ports) & Allocated at run time by \sys{GateTunnel} \\\hline
8100--8199 & Reserved for future framework ranges & Roadmap WG \\\hline
8200--8299 & Inter-cloud session listeners (\sys{GateTunnel}, one per negotiated consumer) & Allocated at run time by \sys{GateTunnel} \\\hline
8300--8429 & Reserved for future framework ranges & Roadmap WG \\\hline
8430--8439 & Life cycle and engineering systems & Roadmap WG \\\hline
8440--8449 & Core systems and trust infrastructure & Roadmap WG \\\hline
8450--8469 & Support systems & Roadmap WG \\\hline
8470--8479 & Monitoring, QoS and analytics systems & Roadmap WG \\\hline
8480--8489 & Health and sidecar endpoints of framework systems & Roadmap WG \\\hline
8490--8499 & Operator tools (LocalCloudManager, LocalCloudConfigurator, SoSConfigurator) & Roadmap WG \\\hline
8500--8799 & Application systems of the local cloud & Local cloud operator \\\hline
8800--8899 & Messaging infrastructure (8883 MQTT over TLS) & Roadmap WG \\\hline
(5671) & AMQP broker for \svc{sysMonitor} (well-known port, outside the block) & Roadmap WG \\\hline
8900--8999 & Reserved & Roadmap WG \\\hline
\end{tabularx}
\caption{Functional ranges of the port block of a local cloud ($k=0$).}
\label{tab:portranges}
\end{table}

Databases and other private back ends of a system are not part of the block: they \emph{should}
not be published outside the host or container network. If they must be, the operator chooses
ports outside the block.

\subsection{Canonical Port Allocation}

{\footnotesize
\begin{longtable}{|P{2.4cm}|P{4.8cm}|P{2.3cm}|P{6.0cm}|}\hline
\rowcolor{gray!33}Port & System & Status & Remark \\\hline\endhead
8432 & \sys{PlantDescriptionEngine} & allocated & as in 4.6 \\\hline
8435 & \sys{Onboarding} & allocated & 4.6 Onboarding controller port; was 8465 in the 5.2.2 migration \\\hline
8437 & (SystemRegistry) & reserved & 4.6, replaced by \sys{ServiceRegistry} \\\hline
8439 & (DeviceRegistry) & reserved & 4.6, replaced by \sys{ServiceRegistry} \\\hline
8441 & \sys{DynamicServiceOrchestration} & allocated & 4.6 Orchestrator port \\\hline
8442 & \sys{SimpleStoreServiceOrchestration} & allocated & was 8465 up to 5.2.2 \\\hline
8443 & \sys{ServiceRegistry} & allocated & bootstrap port \\\hline
8444 & \sys{Authentication} & allocated & \\\hline
8445 & \sys{ConsumerAuthorization} & allocated & 4.6 Authorization port \\\hline
8446 & -- & free (core) & next Core system \\\hline
8447 & (hawkBit configuration manager) & reserved & 4.6 \\\hline
8448 & \sys{CertificateAuthority} & allocated & as in 4.6 \\\hline
8449 & \sys{Gatekeeper} & allocated & as in 4.6 \\\hline
8451 & (QoS Monitor) & reserved & 4.6, replaced by \sys{DeviceQoSEvaluator} \\\hline
8453 & \sys{GateTunnel} & allocated & 4.6 Gateway port \\\hline
8455 & \sys{EventHandler} & allocated & as in 4.6 \\\hline
8457 & \sys{Choreographer} & allocated & as in 4.6 \\\hline
8459 & \sys{Configuration} & allocated & as in 4.6 \\\hline
8461 & \sys{DataManager} & allocated & as in 4.6 \\\hline
8462 & (Translator) & reserved & 4.6, replaced by \sys{TranslationManager} \\\hline
8463 & \sys{TimeManager} & allocated & as in 4.6 \\\hline
8464 & \sys{Blacklist} & allocated & \\\hline
8465 & \sys{TranslationManager} & allocated & sole owner from 5.2.2 on \\\hline
8450, 8452, 8454, 8456, 8458, 8460, 8466--8469 & -- & free (support) & \\\hline
8472 & \sys{DeviceQoSEvaluator} & allocated & \\\hline
8475 & \sys{SysMonitor} & allocated & \\\hline
8478 & \sys{AIGym} & allocated & \\\hline
8470, 8471, 8473, 8474, 8476, 8477, 8479 & -- & free (monitoring) & \\\hline
8480 & \sys{Gatekeeper} health & allocated & not used by the Java-Spring reference (container health checks use the service port) \\\hline
8481 & \sys{GateTunnel} health & allocated & not used by the Java-Spring reference (container health checks use the service port) \\\hline
8482 & \sys{SysMonitor} health & allocated & not used by the Java-Spring reference (container health checks use the service port) \\\hline
8483--8489 & -- & free (health) & \\\hline
8490 & LocalCloudManager UI & allocated & \\\hline
8491 & LocalCloudManager API & allocated & \\\hline
8492 & \sys{LocalCloudConfigurator} & allocated & \\\hline
8493 & \sys{LocalCloudConfigurator} host worker & allocated & loopback of the cloud host only \\\hline
8494 & \sys{SoSConfigurator} & allocated & \\\hline
8495 & \sys{SoSConfigurator} console & allocated & loopback only \\\hline
8496--8499 & -- & free (operator tools) & \\\hline
8501, 8502 & (MSCV, GAMS) & reserved & 4.6 application-level systems \\\hline
8883 & MQTT broker (TLS) & allocated & IANA \emph{secure-mqtt}; \svc{sysMonitorData} and other intra-cloud MQTT traffic \\\hline
5671 & AMQP broker (TLS) & allocated (exception) & IANA \emph{amqps}; carries only \svc{sysMonitor}; outside the block, offset applies ($5671+1000k$) \\\hline
\caption{Canonical port allocation (cloud index $k=0$).}
\label{tab:portalloc}
\end{longtable}
}

Plain MQTT (1883) and other unencrypted infrastructure ports are permitted only for
development on isolated hosts; they are outside the block and are not subject to the cloud
offset.

\subsection{Systems of Local Clouds}
\label{sec:ports:sos}

When several local clouds share hosts or addresses -- test beds, edge servers hosting more
than one cloud, or a system of local clouds deployed on common infrastructure -- their port
blocks \emph{shall not} overlap. Each such cloud is given a \emph{cloud index} $k$ and uses the
block shifted by $1000\cdot k$:
\[
  \mathit{port}(s,k) \;=\; \mathit{canonical}(s) \;+\; 1000\cdot k,
  \qquad k \in \{0,1,\dots,23\}.
\]

\begin{itemize}
  \item The upper bound $k \le 23$ keeps every port below 32768 (P6). The cloud with $k$ uses
        $8000+1000k$ to $8999+1000k$, e.g.\ $k=1$: 9000--9999, $k=2$: 10000--10999.
  \item Clouds on disjoint hosts and addresses \emph{may} all use $k=0$; the cloud index only
        has to be unique among clouds that share an address.
  \item The cloud index is assigned by the operator of the shared infrastructure and recorded
        with the cloud's identity (cloud name, organisation) in the description of the system
        of local clouds. For the 5.2.2 test beds: \texttt{Arrowhead-test-cloud} $k=0$,
        \texttt{Arrowhead-test-cloud-go} $k=1$.
  \item The offset applies to every range of the block, including the \sys{GateTunnel} control
        pool and session listeners, health and operator ports, and the MQTT broker
        (8883 $\rightarrow$ 9883 for $k=1$). It applies equally to the AMQP broker
        (5671 $\rightarrow$ 6671 for $k=1$).
  \item The \emph{inter-cloud relay} is organisation infrastructure shared by the local clouds and
        belongs to none of their blocks. It listens on 8883 (MQTT over TLS) on its own address;
        where it must be published on a host that also carries local clouds, a port outside their
        blocks is used (the test beds use 8893).
  \item Two deployment styles are conformant:
        \emph{(a) offset ports} -- the systems listen on $\mathit{port}(s,k)$ and register it; or
        \emph{(b) canonical inside, offset outside} -- the systems listen on and register the
        canonical port inside an isolated network (e.g.\ a container network), and only the
        published host ports carry the offset. Style (b) satisfies P5 only for consumers inside
        that network; consumers outside it reach the cloud through the offset ports and
        \emph{shall} be configured with the offset \sys{ServiceRegistry} port.
  \item Inter-cloud contacts (\sys{Gatekeeper} cloud records, relays) \emph{shall} carry the
        explicit port of the neighbour cloud's \sys{Gatekeeper}
        ($8449+1000k$); a cloud \emph{shall not} infer it.
\end{itemize}

\subsection{Allocation Procedure}
\begin{enumerate}
  \item A new framework system obtains its canonical port from the free ports of the matching
        range through the Roadmap WG, which updates Table~\ref{tab:portalloc}.
  \item Its implementation \emph{shall} use that port as the shipped default in all
        configuration artefacts (application properties, code defaults, container images,
        compose files and documentation).
  \item Application systems obtain ports from 8500--8799 from the local cloud operator; these
        allocations are local to the cloud.
\end{enumerate}

\subsection{Conformance of the 5.2.2 Implementations}

\begin{table}[h!]
\centering\footnotesize
\begin{tabularx}{\textwidth}{|P{4.3cm}|P{3.3cm}|Y|}\hline
\rowcolor{gray!33}Item & Status & Action \\\hline
\sys{SimpleStoreServiceOrchestration} default 8465 & fixed in workspace & Changed to 8442 (source default, test clouds, image built from the primary Java tree, which also exposes 8442). Only the published upstream image still carries 8465 and copies it at every container start, so deployments of that image set \prop{SERVER_PORT=8442}. Upstream change proposed. \\\hline
\sys{Onboarding} default 8465 & fixed in workspace & Changed to 8435 (code default, properties, IDD, SysDD, test clouds). \\\hline
\sys{TranslationManager} 8465 & conformant & Sole owner of 8465. \\\hline
Test-cloud random-data-producer on container port 8463 (\sys{TimeManager}'s port) & fixed in test clouds & Moved to 8500 (application range, operator allocation); \sys{TimeManager} is now published on its canonical 8463 instead of 8467. \\\hline
\sys{DataManager} on 8462 in the test clouds (\prop{SERVER_PORT=8462}) & fixed in test clouds & Runs on and registers the canonical 8461; 8462 stays reserved. \\\hline
LocalCloudManager UI on 8080, API on 8790 & fixed in test clouds & Moved to 8490 (UI) and 8491 (API), clearing the tunnel pool. \\\hline
Plain MQTT 1883 as \prop{mqtt.broker.port} default in Core configurations & accepted (development) & Production deployments use 8883. \\\hline
Test cloud $k=1$ (\texttt{Arrowhead-test-cloud-go}) & conformant, style (b) & Host ports carry $+1000$ for every system (e.g.\ 9443, 9490, 9883); systems register canonical ports inside the container network. \\\hline
AMQP broker of test cloud $k=1$ published on 9671 & fixed in test clouds & Now published on 6671 ($5671+1000k$). Moving the AMQP broker into the messaging range (e.g.\ 8871) to remove the exception remains a Roadmap WG option. \\\hline
\sys{GateTunnel} control pool configured as 8000--8100 & fixed in test clouds & Upper bound set to 8099 (compose files and code default). \\\hline
Inter-cloud relay published on host port 8893 & conformant & Relay is outside the cloud blocks. \\\hline
Test cloud $k=2$ (\prop{Arrowhead-test-cloud-rust}) on host ports 85xx with container ports 8465 (\sys{SimpleStoreServiceOrchestration}), 8476 (\sys{Onboarding}) and 8480 (LocalCloudManager) & fixed in test clouds & Systems listen on canonical ports (8442, 8435, 8490) and are published on canonical $+2000$ (e.g.\ 10443, UI 10490, API 10491). \\\hline
\end{tabularx}
\caption{Conformance of the 5.2.2 implementations and test beds with the port specification.}
\label{tab:portconf}
\end{table}
